\documentclass[12pt]{article}

% Layout.
\usepackage[top=1in, bottom=0.75in, left=0.75in, right=0.75in, headheight=1in, headsep=6pt]{geometry}

% Fonts.
\usepackage{mathptmx}

\usepackage[scaled=0.86]{helvet}
\renewcommand{\emph}[1]{\textsf{\textbf{#1}}}

% Misc packages.
\usepackage{amsmath,amssymb,latexsym,amsthm}
\usepackage{graphicx,hyperref}
\usepackage{array}
\usepackage{xcolor}
\usepackage{multicol}
\usepackage{tabularx,colortbl}
\usepackage{enumitem}
\usepackage{soul,multicol}
\usepackage{setspace}


\hypersetup{
    colorlinks=true,
    linkcolor=blue,
    filecolor=magenta,      
    urlcolor=blue,
    pdftitle={Proofs Worksheet},
    pdfpagemode=FullScreen,
    }

\def\mailto#1{\href{mailto:#1}{#1}}

%% special math symbols
\newcommand{\N}{\mathbb{N}}
\newcommand{\Z}{\mathbb{Z}}
\newcommand{\R}{\mathbb{R}}
\newcommand{\Q}{\mathbb{Q}}
\newcommand{\sse}{\subseteq}
\newcommand{\mcp}{\mathcal{P}}
\newcommand{\ol}[1]{\overline{#1}}

%other specials
\newcommand{\be}{\begin{enumerate}}
\newcommand{\ee}{\end{enumerate}}
\newcommand{\se}{\subseteq}
\newcommand{\spe}{\supseteq}
\newcommand{\ds}{\displaystyle}
\newcommand{\lcm}{\text{lcm}}
%\newcommand{\gcd}{\text{gcd}}

% Paragraph spacing
\parindent 0pt
\parskip 6pt plus 1pt
\def\tableindent{\hskip 0.5 in}
\def\ts{\hskip 1.5 em}

%header
\usepackage{fancyhdr}
\pagestyle{fancy} 
\lhead{\large\sf\textbf{MATH 265: Introduction to Mathematical Proofs}}
\rhead{\large\sf\textbf{Notes: Ch 14 \S 1-2}}

\newcommand{\localhead}[1]{\par\smallskip\textbf{#1}\nobreak\\}%
\def\heading#1{\localhead{\large\emph{#1}}}
\def\subheading#1{\localhead{\emph{#1}}}

\newenvironment{clist}%
{\bgroup\parskip 0pt\begin{list}{$\bullet$}{\partopsep 4pt\topsep 0pt\itemsep -2pt}}%
{\end{list}\egroup}%

\begin{document}
\begin{enumerate}
\item Two sets $A$ and $B$ have the same \textbf{cardinality} if\\
\vspace{0.6in}

\item Theorem: $|\N| \neq |\R|.$\\
\begin{proof} \quad \\

\vfill

\begin{tabular}{r|l}
$n$ & $f(n)$ \\
\hline
1 & 0.4000000000000000\ldots \\
2 & 8.5006070866690 0\ldots \\
3 & 7.5050094004410 1\ldots \\
4 & 5.5070400804805 0\ldots \\
5 & 6.9002600000050 6\ldots \\
6 & 6.8280958205002 0\ldots \\
7 & 6.5050555065580 8\ldots \\
8 & 8.7208064000044 8\ldots \\
9 & 0.5500008888007 7\ldots \\
10 & 0.5002072207805 1\ldots \\
11 & 2.9000088000090 0\ldots \\
12 & 6.5028000800967 1\ldots \\
13 & 8.8900802400805 0\ldots \\
14 & 8.5000874208022 6\ldots \\
$\vdots$ & $\vdots$ \\
\end{tabular}


\vfill
\vfill
\end{proof}
\item The set $A$ is called \textbf{countably infinite} if \\
\vfill
\item Examples: 
\newpage
\item Theorem 14.3: A set $A$ is countably infinite if and only if 
\vfill
 \item Theorem 14.4: $|\Q|$ is \\
 
 \begin{proof} \quad \\
 

\begin{tabular}{ccccccccccccc}
$0$ & $1$ & $-1$ & $2$ & $-2$ & $3$ & $-3$ & $4$ & $-4$ & $5$ & $-5$ & $\cdots$ \\
\hline
\\
$\dfrac{0}{1}$ & $\dfrac{1}{1}$ & $\dfrac{-1}{1}$ & $\dfrac{2}{1}$ & $\dfrac{-2}{1}$ & $\dfrac{3}{1}$ & $\dfrac{-3}{1}$ & $\dfrac{4}{1}$ & $\dfrac{-4}{1}$ & $\dfrac{5}{1}$ & $\dfrac{-5}{1}$ & $\cdots$ \\[18pt]
 & $\dfrac{1}{2}$ & $\dfrac{-1}{2}$ & $\dfrac{2}{3}$ & $\dfrac{-2}{3}$ & $\dfrac{3}{2}$ & $\dfrac{-3}{2}$ & $\dfrac{4}{3}$ & $\dfrac{-4}{3}$ & $\dfrac{5}{2}$ & $\dfrac{-5}{2}$ & $\cdots$ \\[18pt]
 & $\dfrac{1}{3}$ & $\dfrac{-1}{3}$ & $\dfrac{2}{5}$ & $\dfrac{-2}{5}$ & $\dfrac{3}{4}$ & $\dfrac{-3}{4}$ & $\dfrac{4}{5}$ & $\dfrac{-4}{5}$ & $\dfrac{5}{3}$ & $\dfrac{-5}{3}$ & $\cdots$ \\[18pt]
 & $\dfrac{1}{4}$ & $\dfrac{-1}{4}$ & $\dfrac{2}{7}$ & $\dfrac{-2}{7}$ & $\dfrac{3}{5}$ & $\dfrac{-3}{5}$ & $\dfrac{4}{7}$ & $\dfrac{-4}{7}$ & $\dfrac{5}{4}$ & $\dfrac{-5}{4}$ & $\cdots$ \\[18pt]
 & $\dfrac{1}{5}$ & $\dfrac{-1}{5}$ & $\dfrac{2}{9}$ & $\dfrac{-2}{9}$ & $\dfrac{3}{7}$ & $\dfrac{-3}{7}$ & $\dfrac{4}{9}$ & $\dfrac{-4}{9}$ & $\dfrac{5}{6}$ & $\dfrac{-5}{6}$ & $\cdots$ \\[18pt]
 & $\dfrac{1}{6}$ & $\dfrac{-1}{6}$ & $\dfrac{2}{11}$ & $\dfrac{-2}{11}$ & $\dfrac{3}{8}$ & $\dfrac{-3}{8}$ & $\dfrac{4}{11}$ & $\dfrac{-4}{11}$ & $\dfrac{5}{7}$ & $\dfrac{-5}{7}$ & $\cdots$ \\[18pt]
 & $\dfrac{1}{7}$ & $\dfrac{-1}{7}$ & $\dfrac{2}{13}$ & $\dfrac{-2}{13}$ & $\dfrac{3}{10}$ & $\dfrac{-3}{10}$ & $\dfrac{4}{13}$ & $\dfrac{-4}{13}$ & $\dfrac{5}{8}$ & $\dfrac{-5}{8}$ & $\cdots$ \\[18pt]
 & $\vdots$ & $\vdots$ & $\vdots$ & $\vdots$ & $\vdots$ & $\vdots$ & $\vdots$ & $\vdots$ & $\vdots$ & $\vdots$ & $\ddots$ \\
\end{tabular}

 \vfill
 \end{proof}
\end{enumerate}
\end{document}










